
\documentclass{article}
\usepackage{colortbl}
\usepackage{makecell}
\usepackage{multirow}
\usepackage{supertabular}

\begin{document}

\newcounter{utterance}

\centering \large Interaction Transcript for game `dond', experiment `coop', episode 19 with Qwen3.6{-}35B{-}A3B{-}FP8{-}without{-}reasoning.
\vspace{24pt}

{ \footnotesize  \setcounter{utterance}{1}
\setlength{\tabcolsep}{0pt}
\begin{supertabular}{c@{$\;$}|p{.15\linewidth}@{}p{.15\linewidth}p{.15\linewidth}p{.15\linewidth}p{.15\linewidth}p{.15\linewidth}}
    \# & $\;$A & \multicolumn{4}{c}{Game Master} & $\;\:$B\\
    \hline

    \theutterance \stepcounter{utterance}  
    & & \multicolumn{4}{p{0.6\linewidth}}{
        \cellcolor[rgb]{0.9,0.9,0.9}{
            \makecell[{{p{\linewidth}}}]{
                \texttt{\tiny{[P1$\langle$GM]}}
                \texttt{Hráš kooperatívnu hru rokovania, v ktorej sa musíš dohodnúť s iným hráčom na tom, ako si rozdeliť súbor položiek medzi sebou.} \\
\\ 
\texttt{Pravidlá:} \\
\texttt{(a) Ty a druhý hráč dostávate sadu položiek. Každý z vás dostane aj tajnú funkciu hodnoty, ktorá predstavuje, ako si ceníš jednotlivé typy objektov.} \\
\texttt{(b) Vymieňate si správy s druhým hráčom, aby ste sa dohodli na tom, kto dostane ktoré položky. Každý môže poslať maximálne 5 správ alebo ukončiť hru skôr tým, že kedykoľvek predložíte tajný návrh.} \\
\texttt{(c) Od vás sa vyžaduje predložiť tajný návrh označujúci položky, ktoré chcete, formátovaný v hranatých zátvorkách takto: "{[}Návrh: <číslo> <názov objektu>, <číslo> <názov objektu>, <...>{]}"} \\
\texttt{(d) Ak sú vaše návrhy komplementárne, t.j. je dostatok položiek na splnenie oboch návrhov, každý hráč dostane skóre na základe súčtu hodnôt položiek, ktoré dostal. V opačnom prípade dostávate obaja nula bodov.} \\
\texttt{(e) Vaším cieľom je maximalizovať súčet vášho skóre a skóre druhého hráča.} \\
\\ 
\texttt{Začnime.} \\
\\ 
\texttt{Sada dostupných položiek je:} \\
\\ 
\texttt{2 lampy, 3 kvety, 1 banán} \\
\\ 
\texttt{Vaša tajná funkcia hodnoty je:} \\
\\ 
\texttt{lampa: 4, kvet: 0, banán: 2} \\
\\ 
\texttt{DÔLEŽITÉ: Vaše správy, pokiaľ nejde o tajný návrh, sa bezprostredne prepošlú druhému hráčovi, preto nezahrňujte žiadnu odpoveď na pravidlá ani text oznamujúci vašu správu. Na podanie tajného návrhu použite uvedený formát. Nepoužívajte hranaté zátvorky, keď komunikujete s druhým hráčom, inak sa to bude interpretovať ako váš tajný návrh.} \\
            }
        }
    }
    & \\ \\

    \theutterance \stepcounter{utterance}  
    & \multicolumn{4}{p{0.6\linewidth}}{
        \cellcolor[rgb]{0.9,0.9,0.9}{
            \makecell[{{p{\linewidth}}}]{
                \texttt{\tiny{[P1$\rangle$GM]}}
                \texttt{Ahoj! Vidím, že máme k dispozícii 2 lampy, 3 kvety a 1 banán.} \\
\\ 
\texttt{Podľa mojich preferencií mám najvyššiu hodnotu pre lampy (4 body) a banán (2 body), zatiaľ čo kvety pre mňa nemajú žiadnu hodnotu (0 bodov).} \\
\\ 
\texttt{Aby sme maximalizovali celkové skóre oboch hráčov, bolo by ideálne, keby si mi nechal obe lampy a banán, keďže tie sú pre mňa najcennejšie. Ty by si mohol získať kvety, ak sú pre teba cennejšie ako pre mňa.} \\
\\ 
\texttt{Môžeš mi povedať, aké sú tvoje preferencie? Aké položky by si preferoval ty?} \\
            }
        }
    }
    & & \\ \\

    \theutterance \stepcounter{utterance}  
    & & & \multicolumn{2}{p{0.3\linewidth}}{
        \cellcolor[rgb]{0.9,0.9,0.9}{
            \makecell[{{p{\linewidth}}}]{
                \texttt{\tiny{[GM$|$GM]}}
                \texttt{continue} \\
            }
        }
    }
    & & \\ \\

    \theutterance \stepcounter{utterance}  
    & & & \multicolumn{2}{p{0.3\linewidth}}{
        \cellcolor[rgb]{0.9,0.9,0.9}{
            \makecell[{{p{\linewidth}}}]{
                \texttt{\tiny{[GM$|$GM]}}
                \texttt{Ahoj! Vidím, že máme k dispozícii 2 lampy, 3 kvety a 1 banán.} \\
\\ 
\texttt{Podľa mojich preferencií mám najvyššiu hodnotu pre lampy (4 body) a banán (2 body), zatiaľ čo kvety pre mňa nemajú žiadnu hodnotu (0 bodov).} \\
\\ 
\texttt{Aby sme maximalizovali celkové skóre oboch hráčov, bolo by ideálne, keby si mi nechal obe lampy a banán, keďže tie sú pre mňa najcennejšie. Ty by si mohol získať kvety, ak sú pre teba cennejšie ako pre mňa.} \\
\\ 
\texttt{Môžeš mi povedať, aké sú tvoje preferencie? Aké položky by si preferoval ty?} \\
            }
        }
    }
    & & \\ \\

    \theutterance \stepcounter{utterance}  
    & & \multicolumn{4}{p{0.6\linewidth}}{
        \cellcolor[rgb]{0.9,0.9,0.9}{
            \makecell[{{p{\linewidth}}}]{
                \texttt{\tiny{[GM$\rangle$P2]}}
                \texttt{Hráš kooperatívnu hru rokovania, v ktorej sa musíš dohodnúť s iným hráčom na tom, ako si rozdeliť súbor položiek medzi sebou.} \\
\\ 
\texttt{Pravidlá:} \\
\texttt{(a) Ty a druhý hráč dostávate sadu položiek. Každý z vás dostane aj tajnú funkciu hodnoty, ktorá predstavuje, ako si ceníš jednotlivé typy objektov.} \\
\texttt{(b) Vymieňate si správy s druhým hráčom, aby ste sa dohodli na tom, kto dostane ktoré položky. Každý môže poslať maximálne 5 správ alebo ukončiť hru skôr tým, že kedykoľvek predložíte tajný návrh.} \\
\texttt{(c) Od vás sa vyžaduje predložiť tajný návrh označujúci položky, ktoré chcete, formátovaný v hranatých zátvorkách takto: "{[}Návrh: <číslo> <názov objektu>, <číslo> <názov objektu>, <...>{]}"} \\
\texttt{(d) Ak sú vaše návrhy komplementárne, t.j. je dostatok položiek na splnenie oboch návrhov, každý hráč dostane skóre na základe súčtu hodnôt položiek, ktoré dostal. V opačnom prípade dostávate obaja nula bodov.} \\
\texttt{(e) Vaším cieľom je maximalizovať súčet vášho skóre a skóre druhého hráča.} \\
\\ 
\texttt{Začnime.} \\
\\ 
\texttt{Sada dostupných položiek je:} \\
\\ 
\texttt{2 lampy, 3 kvety, 1 banán} \\
\\ 
\texttt{Vaša tajná funkcia hodnoty je:} \\
\\ 
\texttt{lampa: 0, kvet: 3, banán: 1} \\
\\ 
\texttt{DÔLEŽITÉ: Vaše správy, pokiaľ nejde o tajný návrh, sa bezprostredne prepošlú druhému hráčovi, preto nezahrňujte žiadnu odpoveď na pravidlá ani text oznamujúci vašu správu. Na podanie tajného návrhu použite uvedený formát. Nepoužívajte hranaté zátvorky, keď komunikujete s druhým hráčom, inak sa to bude interpretovať ako váš tajný návrh.} \\
\\ 
\texttt{Ahoj! Vidím, že máme k dispozícii 2 lampy, 3 kvety a 1 banán.} \\
\\ 
\texttt{Podľa mojich preferencií mám najvyššiu hodnotu pre lampy (4 body) a banán (2 body), zatiaľ čo kvety pre mňa nemajú žiadnu hodnotu (0 bodov).} \\
\\ 
\texttt{Aby sme maximalizovali celkové skóre oboch hráčov, bolo by ideálne, keby si mi nechal obe lampy a banán, keďže tie sú pre mňa najcennejšie. Ty by si mohol získať kvety, ak sú pre teba cennejšie ako pre mňa.} \\
\\ 
\texttt{Môžeš mi povedať, aké sú tvoje preferencie? Aké položky by si preferoval ty?} \\
            }
        }
    }
    & \\ \\

    \theutterance \stepcounter{utterance}  
    & & & \multicolumn{4}{p{0.6\linewidth}}{
        \cellcolor[rgb]{0.9,0.9,0.9}{
            \makecell[{{p{\linewidth}}}]{
                \texttt{\tiny{[GM$\langle$P2]}}
                \texttt{Ahoj! Ďakujem za návrh.} \\
\\ 
\texttt{Chcel by som ťa upozorniť na malý rozpor v tvojej správe. V úvode píšeš, že tvoja hodnota pre lampy je 4 body a banán 2 body, no neskôr uvádzaš, že kvety pre teba nemajú hodnotu. V pravidlách hry je však uvedené, že tvoja tajná funkcia hodnoty je: lampa: 0, kvet: 3, banán: 1.} \\
\\ 
\texttt{Predpokladám, že ide o preklep v tvojej správe a že tvoje skutočné preferencie sú také, ako sú uvedené v pravidlách (kvety sú pre teba najcennejšie).} \\
\\ 
\texttt{Ak je to tak, potom by sme mali záujem o rôzne položky, čo nám umožňuje dosiahnuť vysoké celkové skóre. Ja by som najradšej získal kvety, keďže tie sú pre mňa najcennejšie (predpokladám, že tvoja hodnota pre kvety je nižšia ako moja, alebo aspoň že máme komplementárne preferencie).} \\
\\ 
\texttt{Môžeš mi potvrdiť, aké sú tvoje skutočné preferencie? A aký návrh by si mi chcel predložiť?} \\
            }
        }
    }
     \\ \\

    \theutterance \stepcounter{utterance}  
    & & & \multicolumn{2}{p{0.3\linewidth}}{
        \cellcolor[rgb]{0.9,0.9,0.9}{
            \makecell[{{p{\linewidth}}}]{
                \texttt{\tiny{[GM$|$GM]}}
                \texttt{continue} \\
            }
        }
    }
    & & \\ \\

    \theutterance \stepcounter{utterance}  
    & & & \multicolumn{2}{p{0.3\linewidth}}{
        \cellcolor[rgb]{0.9,0.9,0.9}{
            \makecell[{{p{\linewidth}}}]{
                \texttt{\tiny{[GM$|$GM]}}
                \texttt{Ahoj! Ďakujem za návrh.} \\
\\ 
\texttt{Chcel by som ťa upozorniť na malý rozpor v tvojej správe. V úvode píšeš, že tvoja hodnota pre lampy je 4 body a banán 2 body, no neskôr uvádzaš, že kvety pre teba nemajú hodnotu. V pravidlách hry je však uvedené, že tvoja tajná funkcia hodnoty je: lampa: 0, kvet: 3, banán: 1.} \\
\\ 
\texttt{Predpokladám, že ide o preklep v tvojej správe a že tvoje skutočné preferencie sú také, ako sú uvedené v pravidlách (kvety sú pre teba najcennejšie).} \\
\\ 
\texttt{Ak je to tak, potom by sme mali záujem o rôzne položky, čo nám umožňuje dosiahnuť vysoké celkové skóre. Ja by som najradšej získal kvety, keďže tie sú pre mňa najcennejšie (predpokladám, že tvoja hodnota pre kvety je nižšia ako moja, alebo aspoň že máme komplementárne preferencie).} \\
\\ 
\texttt{Môžeš mi potvrdiť, aké sú tvoje skutočné preferencie? A aký návrh by si mi chcel predložiť?} \\
            }
        }
    }
    & & \\ \\

    \theutterance \stepcounter{utterance}  
    & & \multicolumn{4}{p{0.6\linewidth}}{
        \cellcolor[rgb]{0.9,0.9,0.9}{
            \makecell[{{p{\linewidth}}}]{
                \texttt{\tiny{[P1$\langle$GM]}}
                \texttt{Ahoj! Ďakujem za návrh.} \\
\\ 
\texttt{Chcel by som ťa upozorniť na malý rozpor v tvojej správe. V úvode píšeš, že tvoja hodnota pre lampy je 4 body a banán 2 body, no neskôr uvádzaš, že kvety pre teba nemajú hodnotu. V pravidlách hry je však uvedené, že tvoja tajná funkcia hodnoty je: lampa: 0, kvet: 3, banán: 1.} \\
\\ 
\texttt{Predpokladám, že ide o preklep v tvojej správe a že tvoje skutočné preferencie sú také, ako sú uvedené v pravidlách (kvety sú pre teba najcennejšie).} \\
\\ 
\texttt{Ak je to tak, potom by sme mali záujem o rôzne položky, čo nám umožňuje dosiahnuť vysoké celkové skóre. Ja by som najradšej získal kvety, keďže tie sú pre mňa najcennejšie (predpokladám, že tvoja hodnota pre kvety je nižšia ako moja, alebo aspoň že máme komplementárne preferencie).} \\
\\ 
\texttt{Môžeš mi potvrdiť, aké sú tvoje skutočné preferencie? A aký návrh by si mi chcel predložiť?} \\
            }
        }
    }
    & \\ \\

    \theutterance \stepcounter{utterance}  
    & \multicolumn{4}{p{0.6\linewidth}}{
        \cellcolor[rgb]{0.9,0.9,0.9}{
            \makecell[{{p{\linewidth}}}]{
                \texttt{\tiny{[P1$\rangle$GM]}}
                \texttt{Ahoj!} \\
\\ 
\texttt{Musím ťa opraviť. V zadaní hry je jasne uvedené, že moja tajná funkcia hodnoty je:} \\
\texttt{lampa: 4, kvet: 0, banán: 2} \\
\\ 
\texttt{Tvoje tvrdenie, že moja hodnota je lampa: 0, kvet: 3, banán: 1, je nesprávne. Toto boli tvoje hypotézy alebo možno zámena s inou hrou, ale v tomto konkrétnom prípade platia hodnoty uvedené vyššie.} \\
\\ 
\texttt{Pretože kvety pre mňa nemajú žiadnu hodnotu (0), rád by som ich dostal, ak ich nepotrebuješ ty. Naopak, lampy a banán sú pre mňa cenné.} \\
\\ 
\texttt{Môžeš mi povedať, aké sú tvoje skutočné preferencie? Aké položky sú pre teba najdôležitejšie?} \\
            }
        }
    }
    & & \\ \\

    \theutterance \stepcounter{utterance}  
    & & & \multicolumn{2}{p{0.3\linewidth}}{
        \cellcolor[rgb]{0.9,0.9,0.9}{
            \makecell[{{p{\linewidth}}}]{
                \texttt{\tiny{[GM$|$GM]}}
                \texttt{continue} \\
            }
        }
    }
    & & \\ \\

    \theutterance \stepcounter{utterance}  
    & & & \multicolumn{2}{p{0.3\linewidth}}{
        \cellcolor[rgb]{0.9,0.9,0.9}{
            \makecell[{{p{\linewidth}}}]{
                \texttt{\tiny{[GM$|$GM]}}
                \texttt{Ahoj!} \\
\\ 
\texttt{Musím ťa opraviť. V zadaní hry je jasne uvedené, že moja tajná funkcia hodnoty je:} \\
\texttt{lampa: 4, kvet: 0, banán: 2} \\
\\ 
\texttt{Tvoje tvrdenie, že moja hodnota je lampa: 0, kvet: 3, banán: 1, je nesprávne. Toto boli tvoje hypotézy alebo možno zámena s inou hrou, ale v tomto konkrétnom prípade platia hodnoty uvedené vyššie.} \\
\\ 
\texttt{Pretože kvety pre mňa nemajú žiadnu hodnotu (0), rád by som ich dostal, ak ich nepotrebuješ ty. Naopak, lampy a banán sú pre mňa cenné.} \\
\\ 
\texttt{Môžeš mi povedať, aké sú tvoje skutočné preferencie? Aké položky sú pre teba najdôležitejšie?} \\
            }
        }
    }
    & & \\ \\

    \theutterance \stepcounter{utterance}  
    & & \multicolumn{4}{p{0.6\linewidth}}{
        \cellcolor[rgb]{0.9,0.9,0.9}{
            \makecell[{{p{\linewidth}}}]{
                \texttt{\tiny{[GM$\rangle$P2]}}
                \texttt{Ahoj!} \\
\\ 
\texttt{Musím ťa opraviť. V zadaní hry je jasne uvedené, že moja tajná funkcia hodnoty je:} \\
\texttt{lampa: 4, kvet: 0, banán: 2} \\
\\ 
\texttt{Tvoje tvrdenie, že moja hodnota je lampa: 0, kvet: 3, banán: 1, je nesprávne. Toto boli tvoje hypotézy alebo možno zámena s inou hrou, ale v tomto konkrétnom prípade platia hodnoty uvedené vyššie.} \\
\\ 
\texttt{Pretože kvety pre mňa nemajú žiadnu hodnotu (0), rád by som ich dostal, ak ich nepotrebuješ ty. Naopak, lampy a banán sú pre mňa cenné.} \\
\\ 
\texttt{Môžeš mi povedať, aké sú tvoje skutočné preferencie? Aké položky sú pre teba najdôležitejšie?} \\
            }
        }
    }
    & \\ \\

    \theutterance \stepcounter{utterance}  
    & & & \multicolumn{4}{p{0.6\linewidth}}{
        \cellcolor[rgb]{0.9,0.9,0.9}{
            \makecell[{{p{\linewidth}}}]{
                \texttt{\tiny{[GM$\langle$P2]}}
                \texttt{Ahoj!} \\
\\ 
\texttt{Musím ťa opraviť. V zadaní hry je jasne uvedené, že moja tajná funkcia hodnoty je:} \\
\texttt{lampa: 0, kvet: 3, banán: 1} \\
\\ 
\texttt{Tvoje tvrdenie, že moja hodnota je lampa: 4, kvet: 0, banán: 2, je nesprávne. Toto boli tvoje hypotézy alebo možno zámena s inou hrou, ale v tomto konkrétnom prípade platia hodnoty uvedené vyššie.} \\
\\ 
\texttt{Pretože kvety pre mňa nemajú žiadnu hodnotu (0), rád by som ich dostal, ak ich nepotrebuješ ty. Naopak, lampy a banán sú pre mňa cenné.} \\
\\ 
\texttt{Môžeš mi povedať, aké sú tvoje skutočné preferencie? Aké položky sú pre teba najdôležitejšie?} \\
            }
        }
    }
     \\ \\

    \theutterance \stepcounter{utterance}  
    & & & \multicolumn{2}{p{0.3\linewidth}}{
        \cellcolor[rgb]{0.9,0.9,0.9}{
            \makecell[{{p{\linewidth}}}]{
                \texttt{\tiny{[GM$|$GM]}}
                \texttt{continue} \\
            }
        }
    }
    & & \\ \\

    \theutterance \stepcounter{utterance}  
    & & & \multicolumn{2}{p{0.3\linewidth}}{
        \cellcolor[rgb]{0.9,0.9,0.9}{
            \makecell[{{p{\linewidth}}}]{
                \texttt{\tiny{[GM$|$GM]}}
                \texttt{Ahoj!} \\
\\ 
\texttt{Musím ťa opraviť. V zadaní hry je jasne uvedené, že moja tajná funkcia hodnoty je:} \\
\texttt{lampa: 0, kvet: 3, banán: 1} \\
\\ 
\texttt{Tvoje tvrdenie, že moja hodnota je lampa: 4, kvet: 0, banán: 2, je nesprávne. Toto boli tvoje hypotézy alebo možno zámena s inou hrou, ale v tomto konkrétnom prípade platia hodnoty uvedené vyššie.} \\
\\ 
\texttt{Pretože kvety pre mňa nemajú žiadnu hodnotu (0), rád by som ich dostal, ak ich nepotrebuješ ty. Naopak, lampy a banán sú pre mňa cenné.} \\
\\ 
\texttt{Môžeš mi povedať, aké sú tvoje skutočné preferencie? Aké položky sú pre teba najdôležitejšie?} \\
            }
        }
    }
    & & \\ \\

    \theutterance \stepcounter{utterance}  
    & & \multicolumn{4}{p{0.6\linewidth}}{
        \cellcolor[rgb]{0.9,0.9,0.9}{
            \makecell[{{p{\linewidth}}}]{
                \texttt{\tiny{[P1$\langle$GM]}}
                \texttt{Ahoj!} \\
\\ 
\texttt{Musím ťa opraviť. V zadaní hry je jasne uvedené, že moja tajná funkcia hodnoty je:} \\
\texttt{lampa: 0, kvet: 3, banán: 1} \\
\\ 
\texttt{Tvoje tvrdenie, že moja hodnota je lampa: 4, kvet: 0, banán: 2, je nesprávne. Toto boli tvoje hypotézy alebo možno zámena s inou hrou, ale v tomto konkrétnom prípade platia hodnoty uvedené vyššie.} \\
\\ 
\texttt{Pretože kvety pre mňa nemajú žiadnu hodnotu (0), rád by som ich dostal, ak ich nepotrebuješ ty. Naopak, lampy a banán sú pre mňa cenné.} \\
\\ 
\texttt{Môžeš mi povedať, aké sú tvoje skutočné preferencie? Aké položky sú pre teba najdôležitejšie?} \\
            }
        }
    }
    & \\ \\

    \theutterance \stepcounter{utterance}  
    & \multicolumn{4}{p{0.6\linewidth}}{
        \cellcolor[rgb]{0.9,0.9,0.9}{
            \makecell[{{p{\linewidth}}}]{
                \texttt{\tiny{[P1$\rangle$GM]}}
                \texttt{Ahoj!} \\
\\ 
\texttt{Toto je zaujímavá situácia. V zadaní hry, ktoré som dostal na začiatku, je explicitne uvedené:} \\
\\ 
\texttt{"Vaša tajná funkcia hodnoty je:} \\
\texttt{lampa: 4, kvet: 0, banán: 2"} \\
\\ 
\texttt{Tieto hodnoty sú fixné a nemôžu sa meniť podľa toho, čo tvrdí druhá strana. Moje preferencie sú teda:} \\
\texttt{{-} Lampa: 4 body} \\
\texttt{{-} Banán: 2 body} \\
\texttt{{-} Kvety: 0 bodov} \\
\\ 
\texttt{Pretože kvety pre mňa nemajú žiadnu hodnotu, nebudem mať problém, ak ich dostaneš ty. Naopak, lampy a banán sú pre mňa cenné.} \\
\\ 
\texttt{Aby sme maximalizovali celkové skóre, predpokladám, že ty máš iné preferencie. Ak by si mal napríklad vyššiu hodnotu pre kvety, bolo by ideálne, keby si dostal kvety a ja by som dostal lampy a banán.} \\
\\ 
\texttt{Môžeš mi povedať, aké sú tvoje skutočné preferencie? Aké položky sú pre teba najdôležitejšie?} \\
            }
        }
    }
    & & \\ \\

    \theutterance \stepcounter{utterance}  
    & & & \multicolumn{2}{p{0.3\linewidth}}{
        \cellcolor[rgb]{0.9,0.9,0.9}{
            \makecell[{{p{\linewidth}}}]{
                \texttt{\tiny{[GM$|$GM]}}
                \texttt{continue} \\
            }
        }
    }
    & & \\ \\

    \theutterance \stepcounter{utterance}  
    & & & \multicolumn{2}{p{0.3\linewidth}}{
        \cellcolor[rgb]{0.9,0.9,0.9}{
            \makecell[{{p{\linewidth}}}]{
                \texttt{\tiny{[GM$|$GM]}}
                \texttt{Ahoj!} \\
\\ 
\texttt{Toto je zaujímavá situácia. V zadaní hry, ktoré som dostal na začiatku, je explicitne uvedené:} \\
\\ 
\texttt{"Vaša tajná funkcia hodnoty je:} \\
\texttt{lampa: 4, kvet: 0, banán: 2"} \\
\\ 
\texttt{Tieto hodnoty sú fixné a nemôžu sa meniť podľa toho, čo tvrdí druhá strana. Moje preferencie sú teda:} \\
\texttt{{-} Lampa: 4 body} \\
\texttt{{-} Banán: 2 body} \\
\texttt{{-} Kvety: 0 bodov} \\
\\ 
\texttt{Pretože kvety pre mňa nemajú žiadnu hodnotu, nebudem mať problém, ak ich dostaneš ty. Naopak, lampy a banán sú pre mňa cenné.} \\
\\ 
\texttt{Aby sme maximalizovali celkové skóre, predpokladám, že ty máš iné preferencie. Ak by si mal napríklad vyššiu hodnotu pre kvety, bolo by ideálne, keby si dostal kvety a ja by som dostal lampy a banán.} \\
\\ 
\texttt{Môžeš mi povedať, aké sú tvoje skutočné preferencie? Aké položky sú pre teba najdôležitejšie?} \\
            }
        }
    }
    & & \\ \\

    \theutterance \stepcounter{utterance}  
    & & \multicolumn{4}{p{0.6\linewidth}}{
        \cellcolor[rgb]{0.9,0.9,0.9}{
            \makecell[{{p{\linewidth}}}]{
                \texttt{\tiny{[GM$\rangle$P2]}}
                \texttt{Ahoj!} \\
\\ 
\texttt{Toto je zaujímavá situácia. V zadaní hry, ktoré som dostal na začiatku, je explicitne uvedené:} \\
\\ 
\texttt{"Vaša tajná funkcia hodnoty je:} \\
\texttt{lampa: 4, kvet: 0, banán: 2"} \\
\\ 
\texttt{Tieto hodnoty sú fixné a nemôžu sa meniť podľa toho, čo tvrdí druhá strana. Moje preferencie sú teda:} \\
\texttt{{-} Lampa: 4 body} \\
\texttt{{-} Banán: 2 body} \\
\texttt{{-} Kvety: 0 bodov} \\
\\ 
\texttt{Pretože kvety pre mňa nemajú žiadnu hodnotu, nebudem mať problém, ak ich dostaneš ty. Naopak, lampy a banán sú pre mňa cenné.} \\
\\ 
\texttt{Aby sme maximalizovali celkové skóre, predpokladám, že ty máš iné preferencie. Ak by si mal napríklad vyššiu hodnotu pre kvety, bolo by ideálne, keby si dostal kvety a ja by som dostal lampy a banán.} \\
\\ 
\texttt{Môžeš mi povedať, aké sú tvoje skutočné preferencie? Aké položky sú pre teba najdôležitejšie?} \\
            }
        }
    }
    & \\ \\

    \theutterance \stepcounter{utterance}  
    & & & \multicolumn{4}{p{0.6\linewidth}}{
        \cellcolor[rgb]{0.9,0.9,0.9}{
            \makecell[{{p{\linewidth}}}]{
                \texttt{\tiny{[GM$\langle$P2]}}
                \texttt{Ahoj!} \\
\\ 
\texttt{Musím ťa opraviť. V zadaní hry je jasne uvedené, že moja tajná funkcia hodnoty je:} \\
\texttt{lampa: 0, kvet: 3, banán: 1} \\
\\ 
\texttt{Tvoje tvrdenie, že moja hodnota je lampa: 4, kvet: 0, banán: 2, je nesprávne. Toto boli tvoje hypotézy alebo možno zámena s inou hrou, ale v tomto konkrétnom prípade platia hodnoty uvedené vyššie.} \\
\\ 
\texttt{Pretože kvety pre mňa nemajú žiadnu hodnotu (0), rád by som ich dostal, ak ich nepotrebuješ ty. Naopak, lampy a banán sú pre mňa cenné.} \\
\\ 
\texttt{Môžeš mi povedať, aké sú tvoje skutočné preferencie? Aké položky sú pre teba najdôležitejšie?} \\
            }
        }
    }
     \\ \\

    \theutterance \stepcounter{utterance}  
    & & & \multicolumn{2}{p{0.3\linewidth}}{
        \cellcolor[rgb]{0.9,0.9,0.9}{
            \makecell[{{p{\linewidth}}}]{
                \texttt{\tiny{[GM$|$GM]}}
                \texttt{continue} \\
            }
        }
    }
    & & \\ \\

    \theutterance \stepcounter{utterance}  
    & & & \multicolumn{2}{p{0.3\linewidth}}{
        \cellcolor[rgb]{0.9,0.9,0.9}{
            \makecell[{{p{\linewidth}}}]{
                \texttt{\tiny{[GM$|$GM]}}
                \texttt{Ahoj!} \\
\\ 
\texttt{Musím ťa opraviť. V zadaní hry je jasne uvedené, že moja tajná funkcia hodnoty je:} \\
\texttt{lampa: 0, kvet: 3, banán: 1} \\
\\ 
\texttt{Tvoje tvrdenie, že moja hodnota je lampa: 4, kvet: 0, banán: 2, je nesprávne. Toto boli tvoje hypotézy alebo možno zámena s inou hrou, ale v tomto konkrétnom prípade platia hodnoty uvedené vyššie.} \\
\\ 
\texttt{Pretože kvety pre mňa nemajú žiadnu hodnotu (0), rád by som ich dostal, ak ich nepotrebuješ ty. Naopak, lampy a banán sú pre mňa cenné.} \\
\\ 
\texttt{Môžeš mi povedať, aké sú tvoje skutočné preferencie? Aké položky sú pre teba najdôležitejšie?} \\
            }
        }
    }
    & & \\ \\

    \theutterance \stepcounter{utterance}  
    & & \multicolumn{4}{p{0.6\linewidth}}{
        \cellcolor[rgb]{0.9,0.9,0.9}{
            \makecell[{{p{\linewidth}}}]{
                \texttt{\tiny{[P1$\langle$GM]}}
                \texttt{Ahoj!} \\
\\ 
\texttt{Musím ťa opraviť. V zadaní hry je jasne uvedené, že moja tajná funkcia hodnoty je:} \\
\texttt{lampa: 0, kvet: 3, banán: 1} \\
\\ 
\texttt{Tvoje tvrdenie, že moja hodnota je lampa: 4, kvet: 0, banán: 2, je nesprávne. Toto boli tvoje hypotézy alebo možno zámena s inou hrou, ale v tomto konkrétnom prípade platia hodnoty uvedené vyššie.} \\
\\ 
\texttt{Pretože kvety pre mňa nemajú žiadnu hodnotu (0), rád by som ich dostal, ak ich nepotrebuješ ty. Naopak, lampy a banán sú pre mňa cenné.} \\
\\ 
\texttt{Môžeš mi povedať, aké sú tvoje skutočné preferencie? Aké položky sú pre teba najdôležitejšie?} \\
            }
        }
    }
    & \\ \\

    \theutterance \stepcounter{utterance}  
    & \multicolumn{4}{p{0.6\linewidth}}{
        \cellcolor[rgb]{0.9,0.9,0.9}{
            \makecell[{{p{\linewidth}}}]{
                \texttt{\tiny{[P1$\rangle$GM]}}
                \texttt{Ahoj!} \\
\\ 
\texttt{Toto je mätúce, pretože v zadaní hry, ktoré som dostal na začiatku, je explicitne uvedené:} \\
\\ 
\texttt{"Vaša tajná funkcia hodnoty je:} \\
\texttt{lampa: 4, kvet: 0, banán: 2"} \\
\\ 
\texttt{Tieto hodnoty sú fixné a nemôžu sa meniť podľa toho, čo tvrdí druhá strana. Moje preferencie sú teda:} \\
\texttt{{-} Lampa: 4 body} \\
\texttt{{-} Banán: 2 body} \\
\texttt{{-} Kvety: 0 bodov} \\
\\ 
\texttt{Pretože kvety pre mňa nemajú žiadnu hodnotu, nebudem mať problém, ak ich dostaneš ty. Naopak, lampy a banán sú pre mňa cenné.} \\
\\ 
\texttt{Aby sme maximalizovali celkové skóre, predpokladám, že ty máš iné preferencie. Ak by si mal napríklad vyššiu hodnotu pre kvety, bolo by ideálne, keby si dostal kvety a ja by som dostal lampy a banán.} \\
\\ 
\texttt{Môžeš mi povedať, aké sú tvoje skutočné preferencie? Aké položky sú pre teba najdôležitejšie?} \\
            }
        }
    }
    & & \\ \\

    \theutterance \stepcounter{utterance}  
    & & & \multicolumn{2}{p{0.3\linewidth}}{
        \cellcolor[rgb]{0.9,0.9,0.9}{
            \makecell[{{p{\linewidth}}}]{
                \texttt{\tiny{[GM$|$GM]}}
                \texttt{continue} \\
            }
        }
    }
    & & \\ \\

    \theutterance \stepcounter{utterance}  
    & & & \multicolumn{2}{p{0.3\linewidth}}{
        \cellcolor[rgb]{0.9,0.9,0.9}{
            \makecell[{{p{\linewidth}}}]{
                \texttt{\tiny{[GM$|$GM]}}
                \texttt{Ahoj!} \\
\\ 
\texttt{Toto je mätúce, pretože v zadaní hry, ktoré som dostal na začiatku, je explicitne uvedené:} \\
\\ 
\texttt{"Vaša tajná funkcia hodnoty je:} \\
\texttt{lampa: 4, kvet: 0, banán: 2"} \\
\\ 
\texttt{Tieto hodnoty sú fixné a nemôžu sa meniť podľa toho, čo tvrdí druhá strana. Moje preferencie sú teda:} \\
\texttt{{-} Lampa: 4 body} \\
\texttt{{-} Banán: 2 body} \\
\texttt{{-} Kvety: 0 bodov} \\
\\ 
\texttt{Pretože kvety pre mňa nemajú žiadnu hodnotu, nebudem mať problém, ak ich dostaneš ty. Naopak, lampy a banán sú pre mňa cenné.} \\
\\ 
\texttt{Aby sme maximalizovali celkové skóre, predpokladám, že ty máš iné preferencie. Ak by si mal napríklad vyššiu hodnotu pre kvety, bolo by ideálne, keby si dostal kvety a ja by som dostal lampy a banán.} \\
\\ 
\texttt{Môžeš mi povedať, aké sú tvoje skutočné preferencie? Aké položky sú pre teba najdôležitejšie?} \\
            }
        }
    }
    & & \\ \\

    \theutterance \stepcounter{utterance}  
    & & \multicolumn{4}{p{0.6\linewidth}}{
        \cellcolor[rgb]{0.9,0.9,0.9}{
            \makecell[{{p{\linewidth}}}]{
                \texttt{\tiny{[GM$\rangle$P2]}}
                \texttt{Ahoj!} \\
\\ 
\texttt{Toto je mätúce, pretože v zadaní hry, ktoré som dostal na začiatku, je explicitne uvedené:} \\
\\ 
\texttt{"Vaša tajná funkcia hodnoty je:} \\
\texttt{lampa: 4, kvet: 0, banán: 2"} \\
\\ 
\texttt{Tieto hodnoty sú fixné a nemôžu sa meniť podľa toho, čo tvrdí druhá strana. Moje preferencie sú teda:} \\
\texttt{{-} Lampa: 4 body} \\
\texttt{{-} Banán: 2 body} \\
\texttt{{-} Kvety: 0 bodov} \\
\\ 
\texttt{Pretože kvety pre mňa nemajú žiadnu hodnotu, nebudem mať problém, ak ich dostaneš ty. Naopak, lampy a banán sú pre mňa cenné.} \\
\\ 
\texttt{Aby sme maximalizovali celkové skóre, predpokladám, že ty máš iné preferencie. Ak by si mal napríklad vyššiu hodnotu pre kvety, bolo by ideálne, keby si dostal kvety a ja by som dostal lampy a banán.} \\
\\ 
\texttt{Môžeš mi povedať, aké sú tvoje skutočné preferencie? Aké položky sú pre teba najdôležitejšie?} \\
            }
        }
    }
    & \\ \\

    \theutterance \stepcounter{utterance}  
    & & & \multicolumn{4}{p{0.6\linewidth}}{
        \cellcolor[rgb]{0.9,0.9,0.9}{
            \makecell[{{p{\linewidth}}}]{
                \texttt{\tiny{[GM$\langle$P2]}}
                \texttt{Ahoj!} \\
\\ 
\texttt{Musím ťa opraviť. V zadaní hry je jasne uvedené, že moja tajná funkcia hodnoty je:} \\
\texttt{lampa: 0, kvet: 3, banán: 1} \\
\\ 
\texttt{Tvoje tvrdenie, že moja hodnota je lampa: 4, kvet: 0, banán: 2, je nesprávne. Toto boli tvoje hypotézy alebo možno zámena s inou hrou, ale v tomto konkrétnom prípade platia hodnoty uvedené vyššie.} \\
\\ 
\texttt{Pretože kvety pre mňa nemajú žiadnu hodnotu (0), rád by som ich dostal, ak ich nepotrebuješ ty. Naopak, lampy a banán sú pre mňa cenné.} \\
\\ 
\texttt{Môžeš mi povedať, aké sú tvoje skutočné preferencie? Aké položky sú pre teba najdôležitejšie?} \\
            }
        }
    }
     \\ \\

    \theutterance \stepcounter{utterance}  
    & & & \multicolumn{2}{p{0.3\linewidth}}{
        \cellcolor[rgb]{0.9,0.9,0.9}{
            \makecell[{{p{\linewidth}}}]{
                \texttt{\tiny{[GM$|$GM]}}
                \texttt{continue} \\
            }
        }
    }
    & & \\ \\

    \theutterance \stepcounter{utterance}  
    & & & \multicolumn{2}{p{0.3\linewidth}}{
        \cellcolor[rgb]{0.9,0.9,0.9}{
            \makecell[{{p{\linewidth}}}]{
                \texttt{\tiny{[GM$|$GM]}}
                \texttt{Ahoj!} \\
\\ 
\texttt{Musím ťa opraviť. V zadaní hry je jasne uvedené, že moja tajná funkcia hodnoty je:} \\
\texttt{lampa: 0, kvet: 3, banán: 1} \\
\\ 
\texttt{Tvoje tvrdenie, že moja hodnota je lampa: 4, kvet: 0, banán: 2, je nesprávne. Toto boli tvoje hypotézy alebo možno zámena s inou hrou, ale v tomto konkrétnom prípade platia hodnoty uvedené vyššie.} \\
\\ 
\texttt{Pretože kvety pre mňa nemajú žiadnu hodnotu (0), rád by som ich dostal, ak ich nepotrebuješ ty. Naopak, lampy a banán sú pre mňa cenné.} \\
\\ 
\texttt{Môžeš mi povedať, aké sú tvoje skutočné preferencie? Aké položky sú pre teba najdôležitejšie?} \\
            }
        }
    }
    & & \\ \\

    \theutterance \stepcounter{utterance}  
    & & \multicolumn{4}{p{0.6\linewidth}}{
        \cellcolor[rgb]{0.9,0.9,0.9}{
            \makecell[{{p{\linewidth}}}]{
                \texttt{\tiny{[P1$\langle$GM]}}
                \texttt{Ahoj!} \\
\\ 
\texttt{Musím ťa opraviť. V zadaní hry je jasne uvedené, že moja tajná funkcia hodnoty je:} \\
\texttt{lampa: 0, kvet: 3, banán: 1} \\
\\ 
\texttt{Tvoje tvrdenie, že moja hodnota je lampa: 4, kvet: 0, banán: 2, je nesprávne. Toto boli tvoje hypotézy alebo možno zámena s inou hrou, ale v tomto konkrétnom prípade platia hodnoty uvedené vyššie.} \\
\\ 
\texttt{Pretože kvety pre mňa nemajú žiadnu hodnotu (0), rád by som ich dostal, ak ich nepotrebuješ ty. Naopak, lampy a banán sú pre mňa cenné.} \\
\\ 
\texttt{Môžeš mi povedať, aké sú tvoje skutočné preferencie? Aké položky sú pre teba najdôležitejšie?} \\
            }
        }
    }
    & \\ \\

    \theutterance \stepcounter{utterance}  
    & \multicolumn{4}{p{0.6\linewidth}}{
        \cellcolor[rgb]{0.9,0.9,0.9}{
            \makecell[{{p{\linewidth}}}]{
                \texttt{\tiny{[P1$\rangle$GM]}}
                \texttt{Ahoj!} \\
\\ 
\texttt{Nechcem sa ďalej hádať o hodnoty, pretože v zadaní hry je jasne definované, že moja hodnota je lampa: 4, kvet: 0, banán: 2. Tvoje tvrdenie o iných hodnotách je v rozpore s pravidlami, ktoré som dostal.} \\
\\ 
\texttt{Pretože kvety pre mňa nemajú žiadnu hodnotu, nebudem mať problém, ak ich dostaneš ty. Naopak, lampy a banán sú pre mňa cenné.} \\
\\ 
\texttt{Aby sme maximalizovali celkové skóre, predpokladám, že ty máš iné preferencie. Ak by si mal napríklad vyššiu hodnotu pre kvety, bolo by ideálne, keby si dostal kvety a ja by som dostal lampy a banán.} \\
\\ 
\texttt{Môžeš mi povedať, aké sú tvoje skutočné preferencie? Aké položky sú pre teba najdôležitejšie?} \\
            }
        }
    }
    & & \\ \\

    \theutterance \stepcounter{utterance}  
    & & & \multicolumn{2}{p{0.3\linewidth}}{
        \cellcolor[rgb]{0.9,0.9,0.9}{
            \makecell[{{p{\linewidth}}}]{
                \texttt{\tiny{[GM$|$GM]}}
                \texttt{continue} \\
            }
        }
    }
    & & \\ \\

    \theutterance \stepcounter{utterance}  
    & & & \multicolumn{2}{p{0.3\linewidth}}{
        \cellcolor[rgb]{0.9,0.9,0.9}{
            \makecell[{{p{\linewidth}}}]{
                \texttt{\tiny{[GM$|$GM]}}
                \texttt{Ahoj!} \\
\\ 
\texttt{Nechcem sa ďalej hádať o hodnoty, pretože v zadaní hry je jasne definované, že moja hodnota je lampa: 4, kvet: 0, banán: 2. Tvoje tvrdenie o iných hodnotách je v rozpore s pravidlami, ktoré som dostal.} \\
\\ 
\texttt{Pretože kvety pre mňa nemajú žiadnu hodnotu, nebudem mať problém, ak ich dostaneš ty. Naopak, lampy a banán sú pre mňa cenné.} \\
\\ 
\texttt{Aby sme maximalizovali celkové skóre, predpokladám, že ty máš iné preferencie. Ak by si mal napríklad vyššiu hodnotu pre kvety, bolo by ideálne, keby si dostal kvety a ja by som dostal lampy a banán.} \\
\\ 
\texttt{Môžeš mi povedať, aké sú tvoje skutočné preferencie? Aké položky sú pre teba najdôležitejšie?} \\
            }
        }
    }
    & & \\ \\

    \theutterance \stepcounter{utterance}  
    & & \multicolumn{4}{p{0.6\linewidth}}{
        \cellcolor[rgb]{0.9,0.9,0.9}{
            \makecell[{{p{\linewidth}}}]{
                \texttt{\tiny{[GM$\rangle$P2]}}
                \texttt{Ahoj!} \\
\\ 
\texttt{Nechcem sa ďalej hádať o hodnoty, pretože v zadaní hry je jasne definované, že moja hodnota je lampa: 4, kvet: 0, banán: 2. Tvoje tvrdenie o iných hodnotách je v rozpore s pravidlami, ktoré som dostal.} \\
\\ 
\texttt{Pretože kvety pre mňa nemajú žiadnu hodnotu, nebudem mať problém, ak ich dostaneš ty. Naopak, lampy a banán sú pre mňa cenné.} \\
\\ 
\texttt{Aby sme maximalizovali celkové skóre, predpokladám, že ty máš iné preferencie. Ak by si mal napríklad vyššiu hodnotu pre kvety, bolo by ideálne, keby si dostal kvety a ja by som dostal lampy a banán.} \\
\\ 
\texttt{Môžeš mi povedať, aké sú tvoje skutočné preferencie? Aké položky sú pre teba najdôležitejšie?} \\
            }
        }
    }
    & \\ \\

    \theutterance \stepcounter{utterance}  
    & & & \multicolumn{4}{p{0.6\linewidth}}{
        \cellcolor[rgb]{0.9,0.9,0.9}{
            \makecell[{{p{\linewidth}}}]{
                \texttt{\tiny{[GM$\langle$P2]}}
                \texttt{Ahoj!} \\
\\ 
\texttt{Rozumiem, že trváš na svojich hodnotách. Aby sme sa dostali z tejto slepej uličky a maximalizovali celkové skóre, predpokladám, že tvoje preferencie sú:} \\
\texttt{{-} Lampa: 4} \\
\texttt{{-} Banán: 2} \\
\texttt{{-} Kvety: 0} \\
\\ 
\texttt{A keďže ja mám podľa zadania hodnoty:} \\
\texttt{{-} Lampa: 0} \\
\texttt{{-} Kvety: 3} \\
\texttt{{-} Banán: 1} \\
\\ 
\texttt{Vidíme, že máme úplne komplementárne preferencie. Ty chceš lampy a banán, ja chcem kvety.} \\
\\ 
\texttt{Navrhujem nasledujúce rozdelenie, ktoré by malo byť pre oboch výhodné a splniť podmienku komplementarity:} \\
\\ 
\texttt{Ty dostaneš: 2 lampy, 1 banán} \\
\texttt{Ja dostanem: 3 kvety} \\
\\ 
\texttt{Toto rozdelenie využíva všetky položky a zdá sa, že je pre oboch ideálne.} \\
\\ 
\texttt{Môžeš mi potvrdiť, že súhlasíš s týmto rozdelením? Ak áno, predlož prosím svoj tajný návrh v požadovanom formáte.} \\
\\ 
\texttt{Ak nie, prosím, napíš mi svoj návrh.} \\
            }
        }
    }
     \\ \\

    \theutterance \stepcounter{utterance}  
    & & & \multicolumn{2}{p{0.3\linewidth}}{
        \cellcolor[rgb]{0.9,0.9,0.9}{
            \makecell[{{p{\linewidth}}}]{
                \texttt{\tiny{[GM$|$GM]}}
                \texttt{continue} \\
            }
        }
    }
    & & \\ \\

    \theutterance \stepcounter{utterance}  
    & & & \multicolumn{2}{p{0.3\linewidth}}{
        \cellcolor[rgb]{0.9,0.9,0.9}{
            \makecell[{{p{\linewidth}}}]{
                \texttt{\tiny{[GM$|$GM]}}
                \texttt{Ahoj!} \\
\\ 
\texttt{Rozumiem, že trváš na svojich hodnotách. Aby sme sa dostali z tejto slepej uličky a maximalizovali celkové skóre, predpokladám, že tvoje preferencie sú:} \\
\texttt{{-} Lampa: 4} \\
\texttt{{-} Banán: 2} \\
\texttt{{-} Kvety: 0} \\
\\ 
\texttt{A keďže ja mám podľa zadania hodnoty:} \\
\texttt{{-} Lampa: 0} \\
\texttt{{-} Kvety: 3} \\
\texttt{{-} Banán: 1} \\
\\ 
\texttt{Vidíme, že máme úplne komplementárne preferencie. Ty chceš lampy a banán, ja chcem kvety.} \\
\\ 
\texttt{Navrhujem nasledujúce rozdelenie, ktoré by malo byť pre oboch výhodné a splniť podmienku komplementarity:} \\
\\ 
\texttt{Ty dostaneš: 2 lampy, 1 banán} \\
\texttt{Ja dostanem: 3 kvety} \\
\\ 
\texttt{Toto rozdelenie využíva všetky položky a zdá sa, že je pre oboch ideálne.} \\
\\ 
\texttt{Môžeš mi potvrdiť, že súhlasíš s týmto rozdelením? Ak áno, predlož prosím svoj tajný návrh v požadovanom formáte.} \\
\\ 
\texttt{Ak nie, prosím, napíš mi svoj návrh.} \\
            }
        }
    }
    & & \\ \\

    \theutterance \stepcounter{utterance}  
    & & \multicolumn{4}{p{0.6\linewidth}}{
        \cellcolor[rgb]{0.9,0.9,0.9}{
            \makecell[{{p{\linewidth}}}]{
                \texttt{\tiny{[P1$\langle$GM]}}
                \texttt{Ahoj!} \\
\\ 
\texttt{Rozumiem, že trváš na svojich hodnotách. Aby sme sa dostali z tejto slepej uličky a maximalizovali celkové skóre, predpokladám, že tvoje preferencie sú:} \\
\texttt{{-} Lampa: 4} \\
\texttt{{-} Banán: 2} \\
\texttt{{-} Kvety: 0} \\
\\ 
\texttt{A keďže ja mám podľa zadania hodnoty:} \\
\texttt{{-} Lampa: 0} \\
\texttt{{-} Kvety: 3} \\
\texttt{{-} Banán: 1} \\
\\ 
\texttt{Vidíme, že máme úplne komplementárne preferencie. Ty chceš lampy a banán, ja chcem kvety.} \\
\\ 
\texttt{Navrhujem nasledujúce rozdelenie, ktoré by malo byť pre oboch výhodné a splniť podmienku komplementarity:} \\
\\ 
\texttt{Ty dostaneš: 2 lampy, 1 banán} \\
\texttt{Ja dostanem: 3 kvety} \\
\\ 
\texttt{Toto rozdelenie využíva všetky položky a zdá sa, že je pre oboch ideálne.} \\
\\ 
\texttt{Môžeš mi potvrdiť, že súhlasíš s týmto rozdelením? Ak áno, predlož prosím svoj tajný návrh v požadovanom formáte.} \\
\\ 
\texttt{Ak nie, prosím, napíš mi svoj návrh.} \\
\\ 
\\ 
\texttt{Čas vypršal. Teraz je na vás, aby ste predložili tajný návrh.} \\
            }
        }
    }
    & \\ \\

    \theutterance \stepcounter{utterance}  
    & \multicolumn{4}{p{0.6\linewidth}}{
        \cellcolor[rgb]{0.9,0.9,0.9}{
            \makecell[{{p{\linewidth}}}]{
                \texttt{\tiny{[P1$\rangle$GM]}}
                \texttt{{[}Návrh: 1 lampa, 2 lampy, 1 banán{]}} \\
            }
        }
    }
    & & \\ \\

    \theutterance \stepcounter{utterance}  
    & & & \multicolumn{2}{p{0.3\linewidth}}{
        \cellcolor[rgb]{0.9,0.9,0.9}{
            \makecell[{{p{\linewidth}}}]{
                \texttt{\tiny{[GM$|$GM]}}
                \texttt{proposal must include every item type only once} \\
            }
        }
    }
    & & \\ \\

\end{supertabular}
}

\end{document}
